\begin{problem}[2019年第36届全国中学生物理竞赛复赛][农用平板车简化模型]{105}
农用平板车的简化模型如图a所示，两车轮的半径均为 $r$（忽略内外半径差），质量均为 $m$（车轮辐条的质量可忽略），两轮可绕过其中心的光滑细车轴转动（轴的质量可忽略）；车平板长为 $l$、质量为 $2m$，平板的质心恰好位于车轮的轴上；两车把手（可视为细直杆）的长均为 $2l$、质量均为 $m$，且把手前端与平板对齐．平板、把手和车轴固连成一个整体，车轮、平板和把手各自的质量分布都是均匀的．重力加速度大小为 $g$．
\begin{subproblem}
该平板车的车轮被一装置（图中未画出）卡住而不能前后移动，但仍可绕车轴转动。将把手提至水平位置由静止开始释放，求把手在与水平地面碰撞前的瞬间的转动角速度。
\end{subproblem}
\begin{subproblem}
在把手与水平地面碰撞前的瞬间立即撤去卡住两车轮的装置，同时将车轮和轴锁死，在碰后的瞬间立即解锁，假设碰撞时间较短（但不为零），碰后把手末端在竖直方向不反弹。已知把手与地面、车轮与地面之间的滑动摩擦系数均为 $\mu$ （最大静摩擦力等于滑动摩擦力）。求在车轮从开始运动直至静止的过程中，车轴移动的距离。
\end{subproblem}
\end{problem}